% sdapsclassic: SDAPS document class for questionnaires.
% Options:
%   14pt         — base font size; 14pt is recommended for older adults
%   english      — language used for SDAPS-generated text (e.g. instructions)
%   a4paper      — page size
%   sdaps_style=qr — barcode style printed at the bottom of each page;
%                  qr is a 2-D QR code that encodes page/questionnaire ID
%   checkmode=check — how marks are detected: "check" accepts any mark in the box
\documentclass[
  14pt,
  english,
  a4paper,
  sdaps_style=qr,
  checkmode=check
]{sdapsclassic}

% Charter: a serif font with a large x-height and open letterforms,
% recommended for readability in print questionnaires for older adults.
\usepackage{charter}

% Override the normal font size to 14pt with 18pt leading (line spacing).
% The two numbers are {font size}{baseline skip (line spacing)}.
% Increasing baseline skip beyond 1.2× font size improves readability.
\renewcommand{\normalsize}{\fontsize{14}{18}\selectfont}

\ExplSyntaxOn

% Suppress the "DRAFT" watermark that SDAPS prints before the questionnaire
% has been registered with the sdaps command-line tool.
\bool_gset_false:N \g_sdaps_draft_bool

% Set the size and border thickness of the response circles.
%   width/height — diameter of the circle in mm (default: 3.5mm)
%   linewidth    — thickness of the circle border (default: 1bp)
% singlechoice = radio buttons (select one); multichoice = tick boxes (select many)
\sdaps_context_append:nn { singlechoice } { width=6mm, height=6mm, linewidth=1.5bp }
\sdaps_context_append:nn { multichoice }  { width=6mm, height=6mm, linewidth=1.5bp }

% Set the size of the barcode printed at the bottom of each page.
%   c_sdaps_barcode_height_dim  — height of the bars (default: 6.5mm)
%   c_sdaps_barcode_bar_width_dim — width of each individual bar (default: 0.33mm);
%                                   increasing this makes the whole barcode wider.
%                                   Do not exceed ~0.66mm or it may not fit the margin.
\dim_gset:Nn \c_sdaps_barcode_height_dim  { 13mm }
\dim_gset:Nn \c_sdaps_barcode_bar_width_dim { 0.65mm }

\ExplSyntaxOff

% Tighten vertical spacing to fit the questionnaire on one page.
% \parskip is used by SDAPS as the gap before and after each optionquestion block.
\setlength{\parskip}{2pt}
% Reduce space before/after section headings (KOMA-Script defaults are generous).
% beforeskip: space above the heading; afterskip: space below the heading.
% Subsection spacing matters here because each question label is a \subsection.
\RedeclareSectionCommand[beforeskip=6pt, afterskip=3pt]{section}
\RedeclareSectionCommand[beforeskip=4pt, afterskip=1pt]{subsection}

\title{PACES-S}
\author{}
\date{}

% Suppress automatic section numbering in the output.
\setcounter{secnumdepth}{0}

\begin{document}

% questionnaire: the main SDAPS environment. Every page must be inside this.
%   noinfo — suppresses the "how to fill in this form" instructions box
%            that SDAPS normally prints at the top of the first page.
\begin{questionnaire}[noinfo]

% addinfo: stores key–value metadata in the .sdaps project file.
% scale_id is used by the downstream data pipeline to identify which
% questionnaire this is when multiple scales share a single project.
\addinfo{scale_id}{paces_s}

\section{Staff Use Only}

    % Side-by-side layout for the ID and Week boxes.
    % Each minipage is 48% of the line width; \hfill provides the gap between them.
    \noindent
    \begin{minipage}[t]{0.48\linewidth}
        % \textbox* creates a free-text write-in box for manual entry by staff.
        %   * (star)  — the box does NOT stretch vertically to fill extra space
        %   var=...   — variable name written to the SDAPS data file
        %   {0.8cm}   — height of the write-in box
        %   {ID:}     — label printed to the left of the box
        \textbox*[var=worker_id]{0.8cm}{ID:}
    \end{minipage}\hfill
    \begin{minipage}[t]{0.48\linewidth}
        \textbox*[var=week]{0.8cm}{Week:}
    \end{minipage}

    % Side-by-side layout for the Timepoint label and its radio buttons.
    % The label sits in a 3cm minipage; the options share the remaining width.
    % [c] vertically centers both minipages relative to each other.
    \noindent
    \begin{minipage}[c]{3cm}
        Timepoint:
    \end{minipage}%
    \begin{minipage}[c]{\dimexpr\linewidth-3cm\relax}
        % optionquestion: a single-answer (radio button) question.
        %   cols=3    — arrange the choices in 3 equal columns on one row
        %   var=...   — variable name written to the SDAPS data file
        %   {}        — empty title suppresses the subsection heading above the choices
        \begin{optionquestion}[cols=3, var=timepoint]{}
            % \choiceitem defines one radio-button option.
            % The text is the label printed next to the circle.
            \choiceitem{Pre}
            \choiceitem{Baseline}
            \choiceitem{Post}
        \end{optionquestion}
    \end{minipage}

\section{Scale}
Please rate how you feel at the moment about the activity (music and movement class) you have been doing.

\vspace{1mm}

% Each scale item uses a standalone optionquestion rather than a shared optiongroup,
% because the question text appears as a header above its own row of choices.
% Options:
%   cols=5        — 5 equal columns, one per choice option
%   singlechoice  — radio buttons (select exactly one answer)
%   var=...       — variable name written to the SDAPS data file
%   {question}    — the question text shown as a heading above the choices
%
% \choiceitem{text}        — a standard 1-column choice
% \choicemulticolitem{N}{text} — a choice that spans N columns

\begin{optionquestion}[cols=5, singlechoice, var=paces01]{I enjoy it}
  \choiceitem{Strongly disagree}
  \choiceitem{Disagree}
  \choiceitem{Neither agree\\ nor disagree}
  \choiceitem{Agree}
  \choiceitem{Strongly agree}
\end{optionquestion}

\rule{\linewidth}{1pt}

\begin{optionquestion}[cols=5, singlechoice, var=paces02]{I find it pleasurable}
  \choiceitem{Strongly disagree}
  \choiceitem{Disagree}
  \choiceitem{Neither agree\\ nor disagree}
  \choiceitem{Agree}
  \choiceitem{Strongly agree}
\end{optionquestion}

\rule{\linewidth}{1pt}

\begin{optionquestion}[cols=5, singlechoice, var=paces03]{It is very pleasant}
  \choiceitem{Strongly disagree}
  \choiceitem{Disagree}
  \choiceitem{Neither agree\\ nor disagree}
  \choiceitem{Agree}
  \choiceitem{Strongly agree}
\end{optionquestion}

\rule{\linewidth}{1pt}

\begin{optionquestion}[cols=5, singlechoice, var=paces04]{It feels good}
  \choiceitem{Strongly disagree}
  \choiceitem{Disagree}
  \choiceitem{Neither agree\\ nor disagree}
  \choiceitem{Agree}
  \choiceitem{Strongly agree}
\end{optionquestion}

\end{questionnaire}
\end{document}
